\documentclass[amsmath,amssymb,aps,prd,10pt,twocolumn,showkeys]{revtex4}
\usepackage{graphicx}
\usepackage{mathtools}
\usepackage{verbatim}
\DeclareMathOperator\erfc{erfc}
\DeclareMathOperator\erf{erf}
\DeclareMathOperator{\sgn}{sgn}
\DeclareMathOperator{\snr}{SNR}
\begin{document}

\title{Complex Analysis of Askaryan Radiation: UHECR Reconstruction with Askaryan Radio Array}

\author{Jordan C. Hanson}
\email{jhanson2@whittier.edu}
%\affiliation{Department of Physics and Astronomy, Whittier College}
\author{Damian Iba\~{n}ez-Rodriguez}
\affiliation{Department of Physics and Astronomy, Whittier College}
\date{\today}

\keywords{Ultra-high energy neutrino; Askaryan radiation; Mathematical physics}

\maketitle

\section{Responses for Report of Referee A -- DQ14522/Hanson}

The manuscript presents an analytical model that can describe Askaryan radiation from ultra-high energy cosmic ray (UHECR) cascade in ice and as detected in the Askaryan Radio Array antennas. The work builds upon the previously published references 13, 14 and 15 with additional work, mainly focusing on modeling UHECR radio signal. This work could lead toward successful identification of UHECR, a major source of background in detection of ultra-high energy (UHE) neutrinos using in-ice radio techniques, and hence has significance to the field. I have listed my comments below:

\begin{enumerate}
\item In abstract, ``When UHECR cascades enter solid, RF transparent matter at altitudes where the cascade develops, Askaryan radiation can propagate through the solid matter to RF detectors'' Does it refer only to Askaryan radiation produced by cascade developing in the solid medium or also the Askaryan radiation produced in the air? \\

\textit{Our model refers to Askaryan radiation produced within a medium with constant index of refraction.  In this case, we assume the medium is the firn region of the ice sheet.  This is clarified in the abstract, and in subsequent sections of the manuscript.}

\item In introduction, there is ``...the negative excess charge radiates collectively in the 0.1-1 GHz bandwidth''. In abstract, there is ``radio frequency (RF) pulses in the 0.05-1 GHz bandwidth''. Could author elaborate the difference in the lower bandwidth limit? \\ 

\textit{This is more of a typo than a statement about actual bandwidth.  Technically, the Askaryan spectral amplitude decreases linearly with decreasing frequency, until thermal noise takes over in the 10-50 MHz range.  However, RF noise from the galaxy extends to 80-100 MHz, and detectors like ARA often install high-pass filters with 3 dB frequencies around 100 MHz.  We think 100 MHz, or 0.1 GHz, is the more appropriate number.  This has been changed in a consistent fashion throughout the manuscript.}

\item In introduction, there is ``Due to the unique location of ARA at the South Pole, the contribution from the Askaryan effect dominates the observed RF pulses.'' The explanation is later in Section III (``At 90 degrees S, geomagnetic radiation is expected to be horizontally polarized due to the vertical orientation of the local magnetic field. The vertically polarized (VPol) RF channels are there fore less sensitive to geomagnetic radiation.''). Could author put a hint when it is mentioned first time that the explanation will come up later text or put the explanation closer together?
\item In Fig. 7 right plot, I will suggest putting sigma = 0.6 ns line (thick black line) on top of gray lines (or use different color scheme) so it does not get buried in the plot and is easier to spot. What among peak, width, phase or correlation of the lines (UHECR data line and sigma = 0.6 ns model) also being matched in the plot? The peak to peak amplitude (normalized) between UHECR line and sigma=0.6 line has some offset (~ 10 percent). UHECR data oscillates but sigma=0.6 ns model stays flat in either direction of central peak. Could author add comment on that?
\item In Fig 4, bottom right (ID 2716-58611) does not have reflected radio pulse while other plots in Fig 4 likely have both primary and reflected pulse. Is there a feature in the plot itself that hints at having just a primary pulse or both primary and reflected pulse? Also, does the delay between primary and reflected pulse of around 20 ns have a physical meaning (eg, what is the source of reflection and how far the reflecting interface is from antenna)?
\item In Figure 6 and other similar plots, it may help to zoom in further (eg between 75 and 175 ns) / use different color scheme, to see the match between model and CR event better.
\item In section III, when it says ``... and geomagnetic radiation usually dominates the total electromagnetic pulse when the detector lies far beneath the cascade''.  Could you please point me to which text/plot in reference 3 and 5 is the referred from?
\item Data from HPol RF channels is not included in this analysis. Could the conclusion be different if HPol RF is also included in the analysis? How different is the HPol signals from VPol signals for the event geometry considered in this analysis? Could author add some comments about this in the conclusion?
\end{enumerate}

\section{Responses for Report of Referee B -- DQ14522/Hanson}

The authors compare the experimental time-dependent voltage signals in the ARA dipoles at the South Pole generated by cosmic ray air shower cores that strike the snow surface  to the time dependent waveforms predicted by their analytical model described in a previous paper.

This paper looks to be in an important corroboration of the model, but I do suggest several modifications.

\begin{enumerate}
\item The authors should describe how the ARA waveforms in Fig 3-6 were obtained.  In the current draft, at the very end of the paper, the authors describe that the ARA waveform data was directly obtained from the ARA collaboration.  But I think a reference to a private communication the first time the ARA waveform data is introduced will improve the readability.
\item It may help the reader if the time region of the  ``reflected pulse'' is identified in the waveforms shown in one of the figures in figures 3-6. Also, what is responsible for generating the reflected pulse in the ARA signals? I think the authors did a very nice job of showing that the reflected pulses do not affect the results significantly.
\item The paragraph beginning “Several points…” on page 3, argues that pulser calibration events produce a reliable value for f0, but the authors should elaborate on why the calibration events should produce a reliable estimate of gamma, the decay time of the antenna response.  Why doesn’t the double convolution from a pulser signal produce a value for gamma in this analysis that is too large?  Also, the quoted error on the check is larger than the mean, so the authors should argue why this is a meaningful check.
\item The ARA events are different from what is expected from a neutrino interaction in the ice.  The ARA events are generated by the core of a cosmic ray shower that strikes the surface of the snow.  Therefore, the electron distribution is spread over much larger radius (the Moliere radius in air is larger than in ice),  which effects the bandwidth of frequencies that coherently add, and the shower enters the snow surface and encounters a varying index of refraction through the remainder of the shower development.   Both of these effects will alter the time and angular emission dependence  of the Askaryan electric field.  The authors should comment on these differences between the idealized model of neutrino emission assumed in their earlier papers and emission in this case from air shower cores. Given the differences, might the reader expect that the author’s model should provide even better results for neutrino emission within the ice?
\item Fig 3:  what is the purpose for a smaller insert graphic?  What additional information is revealed?  Also, the figure caption should explain the insert graph, instead of letting the reader try to figure it out.
\item Fig. 3:  The authors should explicitly state in the figure caption that the model amplitude is relative voltage amplitude from equation 7, and the ARA event amplitude is relative voltage amplitude.  (to avoid possible confusion with the electric field).
\item In the conclusion, the authors estimate c=0.3/1.78 m/ns, but the index of refraction n=1.78 is value  for ice at depths below ~100m.  The signals from cosmic ray cores are generated near the surface, where index of refraction is much smaller, closer to n=1.4,  and then increases with depth.  Was the estimate in the paper a typo?  If not, the authors should explain the estimate in greater detail.   There is a similar issue with the estimate for the Cherenkov angle.   My last comment is that the parameter a is in units of meters, but since the snow density near the surface is less than half the ice density (0.92 g/cm3) and varies, it will help the reader
to explain if the estimate of a was for ice density, or the snow near the surface?  Usually, lengths in cascade development is parameterized in terms of kg/m$^2$ to account for varying densities of materials.
\item Typo: Measrure on page 7
\item Overall assessment of paper.  This paper provides the first experimental proof that Askaryan radio pulses generated by air shower cores developing within the first 10m of the snow surface are described by the analytic model developed in earlier papers, and therefore, an important paper to publish.
\end{enumerate}

\end{document}